\scanpage{5}
Par suite, on a
\[
\bigwedge_{\underline O_X}^p J/J^2\simeq f^*\bigl(\bigotimes_i\bigwedge^{\alpha_i} M_{d_i}\bigr)(-\sum\alpha_i d_i)
\]
ou $M_{d_i}$ sont les composantes \unsure{non|des} nulles de $M$, supposées
de degré $\alpha_i$. D'où un isom
\[
\Omega^n_{X/S}\otimes f^*\bigl(\bigotimes_i\bigwedge^{\alpha_i}M_{d_i}\bigr)(-\sum\alpha_i d_i)
\xrightarrow{\sim}f^*\bigl(\bigwedge_{\unsure{\underline O_S|S}}^{r+1}\underline E\bigr)(-r-1)
\]
[\unsure{correspondant à|correspondant :|Correspondant :} un hom qui existe en tous cas dès
qu'on se donne des $M_{d_i}\longrightarrow\mathrm{Symm}^{d_i}(\underline E)$, les $M_{d_i}$ loc.
libres de rang $\alpha_i$, et $X$ défini par \deleted{les}\inserted{l'idéal engendré par les} images
des $M_{d_i}$, posant $p=\sum\alpha_i$, $n=r-p$], i.e. on a
\[
\Omega^n_{X/S}\xrightarrow{\sim}f^*\bigl((\bigwedge_S^{r+1}E)\otimes(\bigotimes_i\bigwedge^{\alpha_i}\check M_{d_i})\bigr)(\sum\alpha_i d_i-r-1)
\]

Dans le cas où
\[
\sum\alpha_i d_i=r+1
\]
on trouve donc
\[
\Omega^{\unsure{n}}_{X/S}\xrightarrow{\sim}f^*(\,-\,)
\]

Supposons p.ex. que $E$ soit le faisceau loc. libre de
rang $3$ associé à une courbe elliptique lisse
sur $S$, donc $r=2$, il n'y a qu'un seul $M_{d_i}$, $d_i=3$,
$\alpha_i=1$, on trouve
\[
\Omega^1_{X/S}\xrightarrow{\sim} f^*(\bigwedge^3 E\otimes\check M)
\]
Or $E$ a une filtration de facteurs $\underline O_S$, $\deleted{\unsure{g|L}}\mathcal L^{\otimes2}$, $\mathcal L^{\otimes3}$,
où $\mathcal L$ est un module inv tel que $f^*(\mathcal L)\simeq\unsure{\mathfrak g|\mathfrak T|\mathscr T|\mathcal T}_{X/S}$, donc
$\bigwedge^3 E\simeq\mathcal L^{\otimes5}$, et comme $\Omega^1_{X/S}\simeq f^*(\check{\mathcal L})$, on aura
$\check{\mathcal L}\xrightarrow{\sim}\mathcal L^{\otimes5}\otimes\check M$ d'où $M\xrightarrow{\sim}\mathcal L^{\otimes6}$.
% archive sheet number: 4

